\begin{abstract}
Cross-layer transcoders (CLTs) make attribution graphs smaller by letting one feature write to every later layer. Prior work on a Boolean toy model showed that this can rewrite a deep multi-hop circuit as a sum of shallow single-hop paths. We ask whether this \emph{depth collapse} grows with the true depth of the computation, and whether it survives a fair comparison. We train attention transformers on Boolean tasks of depth $D=1,\dots,5$ with strict interchange intervention training (SIIT), a bypass constraint and reserved site coordinates. We then \emph{certify} each model's depth with interchange, strictness, bypass and circuit-only necessity tests. In doing so we find that a common 2-hop design (final AND) has nonlinear depth 1, because the linear readout performs the last hop. On \NCertified{} certified models we train per-layer (PLT) and cross-layer transcoders with an identical budget (features per layer, exact TopK $L_0$, steps, data), build Circuit-Tracing feature-to-feature attribution graphs (maximum conservation error \MaxConservationErr{}), and measure an effective graph depth $\Dfeat$. Both replacements reconstruct the MLPs almost perfectly, yet \emph{both} graphs lose depth, and the loss grows with $D$ (DCR $\DCRCLTDFour$ for CLT and $\DCRPLTDFour$ for PLT at $D=4$). Logit-level bypass tests show that the shortcuts are introduced by the replacements, not the model ($r_\mathrm{bypass}$: original $\RbypassAllOrig$, PLT $\RbypassAllPLT$, CLT $\RbypassAllCLT$). Interventional tests separate what the graph claims from what the model does. On GPT-2 small, we retain 10 two-hop and 12 one-hop factual-recall prompts. Depth sensitivity is supported for PLT but not CLT, while causal specificity is supported for CLT but not PLT; graph--intervention correlations are positive but inconclusive. This real-model evaluation is preliminary and uses unequal PLT and CLT parameter budgets. InterpBench was not evaluated.
\keywords{Mechanistic interpretability \and Transcoders \and Attribution graphs \and Causal abstraction}
\end{abstract}
